% !TeX root = closed-systems-from-comparison-completeness.tex
% ============================================================
\chapter*{Introduction}
\addcontentsline{toc}{chapter}{Introduction}
\label{ch:introduction}
% ============================================================
A mathematically closed physical theory must be written without
external scaffolding.
If observers, clocks, laboratories, and comparison protocols are all
internal to the same world being modeled, then no background manifold,
state space linearity, causal order, or field content can be accepted as
primitive by convention alone.
Such structure is admissible only when it is recoverable from
intrinsic relational data.
The external-time and external-reference issues motivating this stance
are familiar from the problem-of-time literature in canonical quantum
gravity \cite{Isham1993,Kuchar1992,Anderson2017}; the present manuscript
uses them only as context for the closure requirement stated below.

This manuscript develops that principle as a classification program.
Starting from binary comparison structure alone, it asks which global
constructions are determined, which are excluded, and where any residual
freedom can still occur.
The central claim is that the admissible architecture is highly rigid:
once closed-world admissibility is imposed, quotient semantics,
transport structure, the first stable obstruction layer, and the
intrinsic--smooth interface are mathematically constrained,
and geometric, field-theoretic, scalar, and realization-level
sectors emerge downstream from those constraints.

\medskip
\noindent
\emph{Guiding question.}
\emph{Given only intrinsic comparison data on a closed world, which
mathematical structures are necessary consequences rather than
auxiliary choices?}
\medskip

\section{Primitive object and admissibility stance}
\begin{definition}[Comparison world]
	\label{def:intro-comparison-world}
	A \emph{comparison world} is a pair
	\[
		(U,\mathcal C),
	\]
	where \(U\neq\varnothing\) is a set of states and
	\(\mathcal C\) is a family of binary comparison predicates
	\[
		c:U\times U\to\{0,1\}.
	\]
	We require \(\mathcal C\) to be \emph{intrinsic}: writing
	\(G=\Aut(U,\mathcal C)\) for the automorphism group of
	\cref{def:intrinsic-symmetry-group}, every predicate is
	\(G\)-invariant,
	\[
		c(gu,gv)=c(u,v)
		\qquad
		\text{for all }c\in\mathcal C,\ g\in G,\ u,v\in U.
	\]
	This is a fixed-point condition: \(G\) is determined by \(\mathcal C\),
	and \(\mathcal C\) is admissible exactly when its predicates are stable
	under the symmetries they generate. A predicate violating it assigns
	comparison values tracking a labeling of states that no symmetry of the
	world preserves; such a predicate is an external evaluator in the sense
	of \textup{(SP1)}, not an intrinsic comparison. Intrinsicality is
	therefore the primitive-level expression of the intrinsicality clause
	\textup{(SP1)}: a comparison is a relation among states, not a
	stipulated labeling of them.
\end{definition}

No topology, differentiable structure, metric, linear structure,
probability measure, or background arena is assumed at this level.
The admissibility stance of the book is strict:
all higher structure must be justified as internally determined by the
comparison system and its closed-world coherence conditions.

\section{Standing principle}
\begin{standingprinciple}[Closed-World Admissibility]
\label{sp:closed-world}
Throughout the manuscript we work under one standing principle, unpacked
into the following clauses.
\begin{enumerate}[label=\textup{(SP\arabic*)}]
\item\label{sp:intrinsic}
      \textbf{Intrinsicality clause.}
      No background manifold, external observer frame, or auxiliary
      state-space primitive is admissible unless reconstructed internally
      from comparison data.
\item\label{sp:compactness}
      \textbf{Compactness clause.}
	  Finite-stage and global admissibility are linked by the exact
	  finite-to-global boundary of the refinement tower. Under
	  intrinsicality this boundary is finitary: by
	  \cref{lem:orbit-separation-finite-stage} every comparison separation is
	  witnessed at a single coordinate, so the refinement-tower limit requires
	  no completion and the finite-to-global boundary of
	  \cref{thm:boundary,thm:stability} holds without separate hypothesis. See
	  \cref{lem:orbit-separation-finite-stage,rem:no-limit-only-separation}.
\item\label{sp:closure}
      \textbf{Comparison-completeness clause.}
	  Closedness is profile-pair completeness (rectangular completeness).
	  Under intrinsicality (\cref{def:intro-comparison-world}) and local
	  distinguishability, this is not an independent hypothesis but a
	  consequence; see \cref{thm:rectangular-completeness-intrinsicality}.
\item\label{sp:quotient}
	      \textbf{Quotient-descent clause.}
      Admissible semantic content is exactly the content descending through
      intrinsic quotient maps.
\item\label{sp:transport}
      \textbf{Transport-visibility clause.}
      Any admissible residue beyond quotient semantics is carried by the
      transport-obstruction hierarchy and its first visible layers.
\end{enumerate}
\end{standingprinciple}
\begin{remark}[How hypotheses are used]
In subsequent chapters, theorem headers use either
\emph{local setup hypotheses} (e.g., choosing objects, indices, or
connected components) or \emph{realization-level technical hypotheses}
(e.g., smooth realizability, jet injectivity). These are not independent
foundational axioms; they are scoped instantiations of
\cref{sp:intrinsic,sp:compactness,sp:closure,sp:quotient,sp:transport}.
Any non-local quantitative boundary condition left after theorem-level
discharge is recorded explicitly in \cref{app:conditional-completion}.
\end{remark}

\section{Methodological strategy}
The derivation is organized as a finite-to-global and algebra-to-geometry
program.
At finite stage, one controls compatibility through Boolean admissibility;
at global stage, one passes through exact compactness boundaries.
On the semantic side, comparison completeness determines canonical
quotient structure.
On the enrichment side, one proves that non-quotient content can occur
through at least one of two loci, possibly both, and no third.
On the transport side, one isolates a canonical filtration and proves
that its first stable visible obstruction is quadratic.
Under compatible smooth realization, the stabilized quadratic carrier
enters the intrinsic--smooth interface: the degree--\(2\)
placement is forced by triangle minimality, commutator filtration, and
Stone-limit stabilization, while faithful realization on that forced
carrier makes nonzero stabilized square classes equivalent to nonzero
realized curvature.  The faithfulness condition itself is then
decomposed and discharged: it is equivalent to the ghost-free regime
consisting of a torsion condition \textup{(T)}, a theorem in the Magnus
regime, and a single intrinsic detectability axiom \textup{(D)}, and
the tower supplies a canonical realization on which the second-jet
comparison map is injective by construction
(\cref{sec:interface-ghosts,sec:interface-detectability,sec:interface-canonical}).
The subsequent macroscopic, scalar, and
realization-level consequences proceed through that realized-channel
package.

At a conceptual level, the logical progression is
\[
		\begin{aligned}
			(U,\mathcal C)
			\;&\Longrightarrow\;\text{admissibility boundary and canonical quotient semantics}
			\\
			\;&\Longrightarrow\;\text{two-locus enrichment and transport filtration}
			\\
			\;&\Longrightarrow\;\text{first stable quadratic obstruction}
			\\
			\;&\Longrightarrow\;\text{intrinsic--smooth interface and macroscopic compatibility laws}
			\\
			\;&\Longrightarrow\;\text{scalar-channel classification, numerical closure, and canonical realization}
			\\
			\;&\Longrightarrow\;\text{state-side observable exhaustivity and the automorphism bridge}.
		\end{aligned}
	\]

\section{Main structural claim}
Theorem~\ref{thm:intro-structural-rigidity} records the structural
content proved across
\cref{ch:foundation,ch:bottom,ch:rotation,ch:closure,ch:locus,ch:lift,ch:triangles,ch:descent_obstruction,ch:bridge,ch:flow,ch:stitching,ch:christoffel,ch:interface,ch:causality,ch:hilbert,ch:einstein,ch:connection,ch:em,ch:matter,ch:matter-numbers,ch:numerical-closure,ch:matter-realization,ch:dynamics-program,app:fusion}.

\begin{theorem}[Structural rigidity of closed relational systems]
	\label{thm:intro-structural-rigidity}
	Let \((U,\mathcal C)\) be a comparison world satisfying the closed-system
	admissibility and comparison-completeness conditions developed in this
	manuscript.
	Then the following structure is determined.
	\begin{enumerate}[label=\textup{(\arabic*)}]
		\item
		      Admissible state content factors through canonical quotient
		      semantics, with intrinsic product decomposition
		      \(U\cong X_A\times X_B\) and physical quotient
		      \(\Phys=(X_A\times X_B)/G\).
		\item
		      Beyond quotient semantics, admissible enrichment occurs through at
		      least one of two loci, possibly both---object-level representative
		      choice and morphism-level transport---with no third locus.
		\item
		      Transport admits a canonical descending filtration
		      \(F^1\supset F^2\supset F^3\supset\cdots\),
		      whose first stable visible obstruction is the quadratic layer
		      \(F^2/F^3\).
		\item
		      The same obstruction appears as finite triangle failure,
		      global loop-descent obstruction, and their bridge identification on
		      minimal subsystems.
		\item
		      The obstruction governs the internal observer-reduction framework and
		      the conditional regime of compatible irreversible descent, and admits
		      a canonical spacetime interleaving realization via a two-parameter
		      inverse-limit object.
		\item
		      Under smooth realization, the stabilized quadratic carrier admits an
		      intrinsic--smooth interface: the refinement tower is proved to
		      stabilize its first nontrivial transport obstruction at degree
		      \(2\), and faithful smooth realization on that forced carrier makes a
		      nonzero stabilized square class equivalent to nonzero realized
		      curvature on the realized degree--\(2\) channel.  The faithfulness
		      condition is itself decomposed into the torsion condition
		      \textup{(T)} and the detectability axiom \textup{(D)}, shown
		      necessary by explicit ghost counter-models, and discharged by a
		      canonical realization constructed from the tower.  This
		      degree--\(2\) realized-channel package feeds causal, Hilbert,
		      Einstein-boundary, connection, electromagnetic, matter-sector, and
		      numerical-closure consequences.
		\item
		      The scalar channel determined by the quadratic carrier is then
		      classified intrinsically, assigned its canonical matter-sector
		      realization role, and finally yields a concrete nonperturbative
		      invariant in the magnetic-confinement appendix.
		\item
		      In the separately marked state-side reconstruction extension,
		      relative to the additional realization data
		      \textup{(R1)}--\textup{(R3)}, the raw limit comparison
		      algebra is proved to be larger than every realized smooth
		      reading. The excess is exactly the fiber-separating sector.
		      Under the scale-coherence condition \textup{(U)}, the maximal
		      readable sector is identified with \(C(M)\), its smooth
		      derivations are identified with vector fields on \(M\), and
		      its transport-stabilizing derivations are identified with the
		      symmetry algebra of the realized connection up to gauge.
	\end{enumerate}
\end{theorem}
\begin{proof}
	Each item is proved in the indicated chapter block of the manuscript.
	Item \textup{(1)} follows from the compactness-to-quotient chain in
	\cref{ch:foundation,ch:bottom,ch:rotation,ch:closure}.
	Item \textup{(2)} is the two-locus classification proved in
	\cref{ch:locus,ch:lift}.
	Items \textup{(3)} and \textup{(4)} are established by the transport
	obstruction chain in
	\cref{ch:triangles,ch:descent_obstruction,ch:bridge}.
		Item \textup{(5)} is derived in
		\cref{ch:flow,ch:stitching}.
		Item \textup{(6)} is proved through the smooth-realization chapter, the
		intrinsic-to-smooth interface, and their macroscopic
		continuations in
		\cref{ch:christoffel,ch:interface,ch:causality,ch:hilbert,ch:einstein,ch:connection,ch:em,ch:matter,ch:matter-numbers,ch:numerical-closure}.
		Item \textup{(7)} is completed by the scalar-classification and
		realization chain in
		\cref{ch:matter,ch:matter-numbers,ch:numerical-closure,ch:matter-realization,app:fusion}.
		Item \textup{(8)} is the state-side reconstruction extension proved in
		\cref{ch:dynamics-program}. It is stated relative to its own standing
		data and does not enter the dependency chain for items
		\textup{(1)}--\textup{(7)}.
		Thus the theorem is exactly the consolidated statement of the chapterwise
		results listed above.
\end{proof}

The point of the theorem is methodological as much as structural:
it does not append geometry, field theory, or matter-sector realization
to an independently chosen background.
It shows how these sectors arise as closure-compatible realizations of
structure already determined at the relational level, with scalar and
realization-level consequences forced downstream of the same
obstruction chain.

\section{What is proved, and what is not assumed}
The manuscript proves theorem-level necessity results under the
explicit admissibility hypotheses developed in the text.
It does not assume the standard ordering
"geometry first, fields second."
Instead, it proves a reverse ordering:
semantic closure and transport obstruction first, then geometric,
field-theoretic, scalar, and realization-level structure, followed by a
separately marked state-side reconstruction extension and appendix-level
applications.

The admissibility hypotheses of this manuscript reduce to a single
primitive condition. By
\cref{lem:orbit-separation-finite-stage,lem:finite-support} and
\cref{thm:rectangular-completeness-intrinsicality}, the compactness
clause \textup{(SP2)} and the comparison-completeness clause
\textup{(SP3)} are consequences of intrinsicality
(\cref{def:intro-comparison-world}) together with local
distinguishability. Intrinsicality is not a hypothesis about the universe
but a condition on what counts as a comparison: that a predicate be
invariant under the symmetries it helps constitute. The program therefore
rests on one identity---to be a distinction is to be a
symmetry-invariant of comparison---which is the intrinsicality clause
\textup{(SP1)} stated as a definition of the primitive rather than a
prohibition on later constructions. The manuscript is not axiom-free; it
has exactly one axiom, and that axiom is the meaning of comparison.

In this sense, the contribution is not only a new derivation of familiar
objects.
It is a sharpened criterion for admissibility in closed-system modeling:
which constructions remain valid when all external reference structures
are removed.

\section{Organization of the manuscript}
\label{sec:paper}
The manuscript is organized as a dependency-driven argument rather than a
collection of independent chapters.

\textbf{Foundational phase.}
\Cref{ch:foundation,ch:bottom,ch:rotation,ch:closure} establish the
finite-to-global admissibility boundary, derive comparison completeness,
force quotient semantics, and isolate the subsystem-attribution boundary.
This phase fixes the semantic backbone of the closed system.

\textbf{Classification and obstruction phase.}
\Cref{ch:locus,ch:lift,ch:triangles,ch:descent_obstruction,ch:bridge}
prove the universal two-locus enrichment classification and identify the
first nontrivial transport obstruction through finite and global
representatives of the same defect.

\textbf{Dynamical and geometric phase.}
\Cref{ch:flow,ch:stitching,ch:christoffel,ch:interface,ch:causality}
translate the obstruction framework into the internal
observer-reduction framework, the conditional regime of compatible
irreversible descent, spacetime interleaving, curvature realization,
and macroscopic causal structure.

\textbf{Macroscopic compatibility and field-theoretic phase.}
\Cref{ch:hilbert,ch:einstein,ch:connection,ch:em} derive determined Hilbert
structure, the Einstein compatibility boundary, connection-first
reconstruction, and electromagnetic phase/charge structure from the same
intrinsic chain.

\textbf{Scalar-sector and numerical-closure phase.}
\Cref{ch:matter,ch:matter-numbers,ch:numerical-closure,ch:matter-realization}
classify primitive matter sectors, derive the quartic depth law and
intrinsic scalar closure, and determine the canonical realization role of
the scalar invariant before any appendix-level application.

\textbf{State-side reconstruction extension.}
\Cref{ch:dynamics-program} studies the limit observable algebra after the
realization stack is in place. It proves raw inexhaustibility, identifies
the fiber-separating excess, isolates the scale-coherence condition
\textup{(U)} selecting the maximal readable sector, and identifies
derivations of the smooth regular sector with vector fields and
transport-stabilizing derivations with connection automorphisms up to
gauge.

\textbf{Concrete downstream realization.}
\Cref{app:fusion} then shows that the scalar channel already fixed by the
main stack yields a concrete nonperturbative adiabatic invariant in
magnetic confinement.

\section{How to read this book}
A natural first reading is theorem-driven:
\cref{ch:foundation,ch:closure,ch:locus,ch:bridge,ch:interface,ch:einstein,ch:em,ch:matter,ch:numerical-closure,ch:dynamics-program}.
This path isolates the main boundary and rigidity statements with minimal
technical detours.

A second reading is construction-driven:
\cref{ch:bottom,ch:lift,ch:descent_obstruction,ch:flow,ch:stitching,ch:christoffel,ch:connection,ch:matter-numbers,ch:matter-realization,ch:dynamics-program,app:fusion}.
This path foregrounds explicit constructions and realization mechanisms.

Both readings are equivalent in conclusion:
for closed relational systems, the admissible global architecture is
determined far more rigidly than conventional background-first
formulations suggest.

\section{Conclusion}
\label{sec:introduction-conclusion}
This introduction states the manuscript's governing theorem program in
exact form: within a closed world, admissible structure must be derived
from intrinsic comparison data and not posited as external background.
The chapter architecture is therefore a proof architecture, ordered from
compactness and quotient semantics through transport obstruction to
geometric, field-theoretic, scalar, realization-level, and state-side
reconstruction results.

Accordingly, \cref{ch:foundation} is logically first: it determines the
finite-to-global admissibility boundary that is reused by every later
construction.
